\seccionanexo{Anexo C: Cuestionarios de Evaluación de la Versión 2.0}

\providecommand{\casilla}{\framebox[0.38cm]{\rule{0pt}{0.33cm}}}

Cuestionarios de preguntas cerradas con diseño de respuesta para medir la satisfacción y el impacto de la Versión 2.0 del sistema web de inscripción. Cada bloque presenta una cuadrícula de casillas que el encuestado marca con una X. Escala Likert de 5 puntos: 1 = Totalmente en desacuerdo, 2 = En desacuerdo, 3 = Neutral, 4 = De acuerdo, 5 = Totalmente de acuerdo.

\Needspace{7cm}
\noindent\textbf{Cuestionario A-1: Evaluación de la Versión 2.0 --- Estudiantes postulantes (n = 60, autoadministrado)}

\begin{table}[H]
    \centering
    \caption{Encabezado e instrucciones del cuestionario A-1 (estudiantes).}
    \label{tab:a1_encabezado}
    \begin{tabularx}{\textwidth}{p{0.26\textwidth} X}
        \toprule
        \textbf{Campo} & \textbf{Contenido} \\
        \midrule
        Institución & MÓDULO TECNOLÓGICO PRODUCTIVO SAN JULIÁN BTH --- Especialidad de Sistemas Informáticos \\
        Instrumento & Cuestionario A-1: Evaluación de la Versión 2.0 del sistema web de inscripción \\
        Muestra & 60 estudiantes postulantes (muestreo por conveniencia) \\
        Aplicación & Primer bimestre de la gestión 2026 \\
        Tiempo estimado & 5 minutos \\
        Instrucciones & Marcar con una X una sola opción por pregunta. Respuestas anónimas (C.I. como ****). \\
        \bottomrule
    \end{tabularx}
\end{table}

\Needspace{12cm}
\noindent\textbf{Bloque A. Identificación (anonimizada):}

\begin{table}[H]
    \centering
    \caption{Diseño de respuesta del bloque A (identificación) del cuestionario A-1.}
    \label{tab:a1_identificacion}
    \renewcommand{\arraystretch}{1.60}
    \begin{tabularx}{\textwidth}{c L}
        \toprule
        \textbf{N.º} & \textbf{Pregunta y opciones de respuesta} \\
        \midrule
        A1 & Unidad educativa de procedencia: \rule[-2pt]{4.5cm}{0.4pt} \\
        \addlinespace
        A2 & Especialidad de interés:\newline
            \mbox{\casilla{} Construcción Civil}\quad \mbox{\casilla{} Contaduría General}\quad \mbox{\casilla{} Electrónica}\quad \mbox{\casilla{} Gastronomía}\quad \mbox{\casilla{} Mecánica Automotriz}\quad \mbox{\casilla{} Sistemas Informáticos}\quad \mbox{\casilla{} Textiles y Confecciones} \\
        \addlinespace
        A3 & Edad:\newline
            \mbox{\casilla{} 15 años o menos}\quad \mbox{\casilla{} 16 años}\quad \mbox{\casilla{} 17 años}\quad \mbox{\casilla{} 18 años o más} \\
        \addlinespace
        A4 & Acceso a internet en su teléfono celular:\newline
            \mbox{\casilla{} Siempre}\quad \mbox{\casilla{} A veces}\quad \mbox{\casilla{} Nunca} \\
        \addlinespace
        A5 & Frecuencia de uso de la página de la V2.0:\newline
            \mbox{\casilla{} 1 vez}\quad \mbox{\casilla{} 2 a 3 veces}\quad \mbox{\casilla{} Más de 3 veces} \\
        \bottomrule
    \end{tabularx}
\end{table}

\Needspace{12cm}
\noindent\textbf{Bloque B. Ítems de evaluación en escala Likert (1--5):}

\begin{table}[H]
    \centering
    \caption{Cuadrícula de respuesta del bloque B (escala Likert 1--5) del cuestionario A-1.}
    \label{tab:a1_likert}
    \renewcommand{\arraystretch}{1.30}
    \small
    \begin{tabularx}{\textwidth}{X c c c c c}
        \toprule
        \textbf{Ítem} & \textbf{1} & \textbf{2} & \textbf{3} & \textbf{4} & \textbf{5} \\
        \midrule
        B1. La interfaz web es clara y se adapta adecuadamente a la pantalla de mi teléfono celular. & \casilla & \casilla & \casilla & \casilla & \casilla \\
        B2. Los pasos para completar el formulario de preinscripción son comprensibles y rápidos de ejecutar. & \casilla & \casilla & \casilla & \casilla & \casilla \\
        B3. Pude consultar mi usuario y restablecer mi contraseña de forma autónoma con mi cédula de identidad y fecha de nacimiento (menos de 2 min). & \casilla & \casilla & \casilla & \casilla & \casilla \\
        B4. El mecanismo de autoservicio me permitió asegurar oportunamente mi cupo en la especialidad técnica de mi preferencia sin hacer filas presenciales. & \casilla & \casilla & \casilla & \casilla & \casilla \\
        B5. El comprobante digital emitido refleja con claridad los requisitos documentales exigidos (libretas de 1ro a 3ro, cédulas, carnet y certificado). & \casilla & \casilla & \casilla & \casilla & \casilla \\
        B6. La entrega y verificación de mis documentos con los tutores de curso se registró de manera ágil y ordenada en el sistema. & \casilla & \casilla & \casilla & \casilla & \casilla \\
        B7. La asignación a mi paralelo (4to A o 4to B) respetó la compatibilidad de horarios con mi colegio de origen. & \casilla & \casilla & \casilla & \casilla & \casilla \\
        B8. La distribución de plazas entre los colegios del distrito municipal fue equitativa y transparente. & \casilla & \casilla & \casilla & \casilla & \casilla \\
        \bottomrule
    \end{tabularx}
    \notatabla{Escala Likert: 1 Totalmente en desacuerdo, 2 En desacuerdo, 3 Neutral, 4 De acuerdo, 5 Totalmente de acuerdo.}
\end{table}

\Needspace{16cm}
\noindent\textbf{Bloque C. Preguntas cerradas de opción única:}

\begin{table}[H]
    \centering
    \caption{Diseño de respuesta del bloque C (preguntas cerradas) del cuestionario A-1.}
    \label{tab:a1_cerradas}
    \renewcommand{\arraystretch}{1.60}
    \begin{tabularx}{\textwidth}{c L}
        \toprule
        \textbf{N.º} & \textbf{Pregunta y opciones de respuesta} \\
        \midrule
        C1 & ¿Cuánto demoró su trámite de inscripción en la página de la V2.0?\newline
            \mbox{\casilla{} Menos de 5 min}\quad \mbox{\casilla{} 5 a 15 min}\quad \mbox{\casilla{} 16 a 30 min}\quad \mbox{\casilla{} Más de 30 min} \\
        \addlinespace
        C2 & ¿Recuperó su contraseña con su cédula de identidad y fecha de nacimiento?\newline
            \mbox{\casilla{} Sí, en menos de 2 min}\quad \mbox{\casilla{} Sí, pero demoró más}\quad \mbox{\casilla{} No lo intenté}\quad \mbox{\casilla{} No pude recuperarla} \\
        \addlinespace
        C3 & En comparación con su inscripción anterior, esta fue:\newline
            \mbox{\casilla{} Mucho más rápida}\quad \mbox{\casilla{} Igual}\quad \mbox{\casilla{} Más lenta}\quad \mbox{\casilla{} No participó antes} \\
        \addlinespace
        C4 & Satisfacción global con la página:\newline
            \mbox{\casilla{} Muy insatisfecho(a)}\quad \mbox{\casilla{} Insatisfecho(a)}\quad \mbox{\casilla{} Neutral}\quad \mbox{\casilla{} Satisfecho(a)}\quad \mbox{\casilla{} Muy satisfecho(a)} \\
        \addlinespace
        C5 & ¿Recomendaría esta página a otros postulantes?\newline
            \mbox{\casilla{} Sí}\quad \mbox{\casilla{} No} \\
        \addlinespace
        C6 & ¿Cuál fue el aspecto más valioso para usted?\newline
            \mbox{\casilla{} Recuperación autónoma de la contraseña}\quad \mbox{\casilla{} Registro y verificación de documentos}\quad \mbox{\casilla{} Asignación automática de paralelos}\quad \mbox{\casilla{} Comprobante con código de verificación} \\
        \bottomrule
    \end{tabularx}
\end{table}

\vspace{0.4cm}
\Needspace{7cm}
\noindent\textbf{Cuestionario A-2: Evaluación de la Versión 2.0 --- Personal administrativo y docente (n = 5, entrevista estandarizada)}

\begin{table}[H]
    \centering
    \caption{Encabezado e instrucciones del cuestionario A-2 (funcionarios).}
    \label{tab:a2_encabezado}
    \begin{tabularx}{\textwidth}{p{0.26\textwidth} X}
        \toprule
        \textbf{Campo} & \textbf{Contenido} \\
        \midrule
        Institución & MÓDULO TECNOLÓGICO PRODUCTIVO SAN JULIÁN BTH --- Especialidad de Sistemas Informáticos \\
        Instrumento & Cuestionario A-2: Evaluación de la Versión 2.0 del sistema web de inscripción \\
        Muestra & 5 funcionarios: Director, Secretaría Académica y 3 docentes (Tutores de Curso) \\
        Aplicación & Primer bimestre de la gestión 2026 \\
        Tiempo estimado & 6 minutos \\
        Instrucciones & El aplicador marca con una X la opción elegida por el entrevistado. Respuestas anónimas (C.I. como ****). \\
        \bottomrule
    \end{tabularx}
\end{table}

\Needspace{6cm}
\noindent\textbf{Bloque A. Identificación (anonimizada):}

\begin{table}[H]
    \centering
    \caption{Diseño de respuesta del bloque A (identificación) del cuestionario A-2.}
    \label{tab:a2_identificacion}
    \renewcommand{\arraystretch}{1.60}
    \begin{tabularx}{\textwidth}{c L}
        \toprule
        \textbf{N.º} & \textbf{Pregunta y opciones de respuesta} \\
        \midrule
        A1 & Cargo:\newline
            \mbox{\casilla{} Director}\quad \mbox{\casilla{} Secretaría Académica}\quad \mbox{\casilla{} Docente (Tutor de Curso)} \\
        \addlinespace
        A2 & Antigüedad en el MTP:\newline
            \mbox{\casilla{} Menos de 5 años}\quad \mbox{\casilla{} 5 a 10 años}\quad \mbox{\casilla{} 11 a 15 años}\quad \mbox{\casilla{} Más de 15 años} \\
        \addlinespace
        A3 & ¿Participó en las pruebas del sistema de la V2.0?\newline
            \mbox{\casilla{} Sí}\quad \mbox{\casilla{} No} \\
        \bottomrule
    \end{tabularx}
\end{table}

\Needspace{11cm}
\noindent\textbf{Bloque B. Ítems de evaluación en escala Likert (1--5):}

\begin{table}[H]
    \centering
    \caption{Cuadrícula de respuesta del bloque B (escala Likert 1--5) del cuestionario A-2.}
    \label{tab:a2_likert}
    \renewcommand{\arraystretch}{1.30}
    \small
    \begin{tabularx}{\textwidth}{X c c c c c}
        \toprule
        \textbf{Ítem} & \textbf{1} & \textbf{2} & \textbf{3} & \textbf{4} & \textbf{5} \\
        \midrule
        B1. El sistema redujo el tiempo de atención presencial en ventanilla para el proceso de inscripción. & \casilla & \casilla & \casilla & \casilla & \casilla \\
        B2. El registro y la verificación de los documentos son más ordenados y comprobables que antes. & \casilla & \casilla & \casilla & \casilla & \casilla \\
        B3. Los reportes permiten detectar expedientes incompletos antes de cerrar el proceso. & \casilla & \casilla & \casilla & \casilla & \casilla \\
        B4. La asignación automática de cursos y paralelos redujo los errores de planilla. & \casilla & \casilla & \casilla & \casilla & \casilla \\
        B5. El autoservicio de contraseñas disminuyó la carga de trabajo de secretaría. & \casilla & \casilla & \casilla & \casilla & \casilla \\
        B6. El sistema es fácil de operar con los recursos disponibles (computadora o celular). & \casilla & \casilla & \casilla & \casilla & \casilla \\
        \bottomrule
    \end{tabularx}
    \notatabla{Escala Likert: 1 Totalmente en desacuerdo, 2 En desacuerdo, 3 Neutral, 4 De acuerdo, 5 Totalmente de acuerdo.}
\end{table}

\Needspace{12cm}
\noindent\textbf{Bloque C. Preguntas cerradas de opción única:}

\begin{table}[H]
    \centering
    \caption{Diseño de respuesta del bloque C (preguntas cerradas) del cuestionario A-2.}
    \label{tab:a2_cerradas}
    \renewcommand{\arraystretch}{1.60}
    \begin{tabularx}{\textwidth}{c L}
        \toprule
        \textbf{N.º} & \textbf{Pregunta y opciones de respuesta} \\
        \midrule
        C1 & ¿Cuántas horas semanales ahorró su área en el proceso de inscripción?\newline
            \mbox{\casilla{} Menos de 2}\quad \mbox{\casilla{} 2 a 4}\quad \mbox{\casilla{} 5 a 8}\quad \mbox{\casilla{} Más de 8} \\
        \addlinespace
        C2 & ¿Qué porcentaje de expedientes quedó con registro completo y verificable?\newline
            \mbox{\casilla{} 0 a 50}\quad \mbox{\casilla{} 51 a 80}\quad \mbox{\casilla{} 81 a 99}\quad \mbox{\casilla{} 100} \\
        \addlinespace
        C3 & ¿Qué errores de asignación de paralelos se observaron desde la V2.0?\newline
            \mbox{\casilla{} Ninguno}\quad \mbox{\casilla{} Pocos}\quad \mbox{\casilla{} Frecuentes}\quad \mbox{\casilla{} Muy frecuentes} \\
        \addlinespace
        C4 & ¿Qué recomienda para el sistema?\newline
            \mbox{\casilla{} Mantenerlo como está}\quad \mbox{\casilla{} Mantenerlo con mejoras}\quad \mbox{\casilla{} Reemplazarlo} \\
        \bottomrule
    \end{tabularx}
\end{table}

\notatabla{Cuestionarios de evaluación de la Versión 2.0 aplicados en el primer bimestre de 2026 a 60 estudiantes postulantes y 5 funcionarios del MTP San Julián (muestreo por conveniencia). Datos anonimizados (C.I. como ****). Validación por juicio del tutor y prueba piloto con 5 estudiantes (ver Anexo B).}
